% ==================================================================
\section{Results and Discussion}
\label{sec:results}
% ==================================================================
%
% Numerical values in Tables III, V (severity), and X (economics) are
% measured from the trained GraphSAGE-MAPPO checkpoint
% (artifacts/checkpoints_am_mappo_v2/am_mappo_final.pt) under the
% AM-only build, using scripts/finalize_paper.py to patch this file from
% paper/figures_real/*.csv. MLP-MAPPO / GCNN-PPO / MATD3 baselines are
% trained separately and integrated in a later revision; the current
% measured comparison covers GraphSAGE-MAPPO + Fixed Droop + No FFR.

This section reports quantitative results from the AM-focused tri-layer
framework. The slow-timescale market clearing dataset is generated in
\texttt{data/precomputed\_365d\_97to67} (365 days, 15-minute resolution,
35\,040 samples; 345 days for training, 20 days held out for evaluation).
The fast-timescale FFR and topology-generalisation tables are produced by
\texttt{src/eval/eval\_ffr\_topology.py} (six methods, four scenarios, 20 runs
per cell). Economic results are produced by \texttt{src/eval/eval\_economics.py}.

\smallskip
\noindent\textbf{Architecture note.} All three layers (L0 DSO pricing,
L1 VPP dispatch, L2 GraphSAGE-MAPPO control) are oriented around the AM
service portfolio (energy P + FFR capacity/activation). A reactive (Q)
product has been intentionally excluded to keep the optimisation
landscape and revenue accounting consistent across layers; reactive
support, when needed, is provided implicitly via inverter Volt/Var
control rather than as a remunerated market product.

% =====================================================================
\subsection{Experimental Setup}
\label{ssec:setup}
% =====================================================================

\textbf{System.} Modified IEEE 123-bus islanded distribution feeder
($S_{\text{BASE}}=15.705$~MW, $H_{\text{SYS}}=1.18$~s, $D_{\text{SYS}}=0.73$~pu,
$R_{\text{SYS}}=0.048$~pu, $f_0=50$~Hz). Five grid-forming (GFM) inverters at
buses 60, 67, 98, 101, 105 aggregated into three VPPs (VPP1--VPP3); 41 RL
agents in total. Per-step inference budget for real-time control: 50~ms.

\textbf{Methods.} Six controllers compared along the ladder of
Table~\ref{tab:methods}: GraphSAGE-MAPPO (proposed), three trained baselines
(MLP-MAPPO ablation, GCNN-PPO~\cite{guo2024gcnn}, MATD3~\cite{li2025matd3}),
and two non-learning controls (fixed 5~\% droop, no FFR).

\begin{table}[!t]
\caption{Method lineup and isolated component}
\label{tab:methods}
\centering
\footnotesize
\begin{tabular}{lllll}
\toprule
\# & Method & Encoder & RL algo. & Multi-agent \\
\midrule
1 & GraphSAGE-MAPPO (\emph{ours}) & GraphSAGE   & PPO     & MAPPO (CTDE) \\
2 & MLP-MAPPO (ablation)          & MLP         & PPO     & MAPPO (CTDE) \\
3 & GCNN-PPO~\cite{guo2024gcnn}   & Spectral GCN& PPO     & Single centralised \\
4 & MATD3~\cite{li2025matd3}      & MLP         & TD3     & CTDE \\
5 & Fixed Droop                   & ---         & ---     & --- \\
6 & No FFR                        & ---         & ---     & --- \\
\bottomrule
\end{tabular}
\end{table}

\textbf{Scenarios.} Four contingencies sized by literature
(Table~\ref{tab:scenarios}): one moderate load step, two generation trips of
increasing severity, and one line trip that mutates the feeder edge index
mid-episode. Magnitudes range from 15--35\,\% of $S_{\text{BASE}}$, all above
the no-FFR analytical RoCoF tolerance of 1.49~MW.

\begin{table}[!t]
\caption{Contingency scenarios}
\label{tab:scenarios}
\centering
\footnotesize
\begin{tabular}{llccc}
\toprule
Sc. & Event       & $\Delta P$ (MW) & \% $S_B$ & Purpose \\
\midrule
S1  & load\_step  & $+2.5$  & 16\,\% & moderate, droop reachable \\
S2  & gen\_trip   & $-3.9$  & 25\,\% & severe; droop fails \\
S3  & line\_trip  & $-2.4$  & 15\,\% & topology mutation \\
S4  & gen\_trip   & $-5.5$  & 35\,\% & extreme; graceful degradation \\
\bottomrule
\end{tabular}
\end{table}

\textbf{Topology split.} 20 reconfigurations from the topology cache are split
via greedy farthest-point selection (15 train, 5 test). Resulting Jaccard
edge-distance statistics are in Table~\ref{tab:split_quality}; the mean
$d_{\min}$ is $3.0\times$ that of a random split, putting test topologies
genuinely out-of-distribution.

\begin{table}[!t]
\caption{Topology split quality}
\label{tab:split_quality}
\centering
\footnotesize
\begin{tabular}{lcccc}
\toprule
Method & $|\mathcal{T}_{\text{train}}|$ & $|\mathcal{T}_{\text{test}}|$ &
$\overline{d_{\min}}$ & $\min\, d_{\min}$ \\
\midrule
Random (legacy)         & 15 & 5 & 0.008 & 0.002 \\
Farthest-point (ours)   & 15 & 5 & \textbf{0.024} & \textbf{0.018} \\
\bottomrule
\end{tabular}
\end{table}

% =====================================================================
\subsection{AM Results: FFR Performance Comparison}
\label{ssec:results_am}
% =====================================================================

Table~\ref{tab:ffr_main} reports the four core frequency-security metrics
per scenario and per method, averaged over 20 runs. The proposed
GraphSAGE-MAPPO is the only method that retains FFR-SR\,=\,100\,\% on the
extreme S4 contingency (5.5\,MW gen-trip): Fixed Droop and No FFR both
drop the nadir below the 49.5\,Hz UFLS bound and score 0\,\% on S4.
Across all four scenarios, the proposed method delivers the lowest
post-event IAE (60--74\,\% reduction vs.\ No FFR on S1/S3, 33--52\,\%
reduction on S2/S4) and reaches the $\pm0.1$\,Hz settling band more than
twice as fast on the mild scenarios (S1, S3). On S2 and S4 the
settling-window estimator saturates at the 50\,s post-window length for
every method, so the cleaner separator there is IAE and nadir margin.
Note that the baseline GFM primary droop is intentionally retained in all
configurations, which is why No FFR and Fixed Droop still achieve nominal
nadir on the lighter contingencies---the differentiation appears in IAE,
settling time, and the extreme-event success rate. MLP-MAPPO, GCNN-PPO,
and MATD3 baselines are pending separate training; their rows will be
added in a subsequent revision.

\begin{table*}[!t]
\caption{FFR performance per scenario and per method (mean over 20 runs, 15 train topologies)}
\label{tab:ffr_main}
\centering
\footnotesize
\begin{tabular}{llcccccc}
\toprule
Scenario & Method & Nadir (Hz) & RoCoF$_{\max}$ (Hz/s) & IAE$_{\text{post}}$ (Hz$\cdot$s) & ITAE (Hz$\cdot$s$^2$) & $t_{\text{settle}}$ (s) & FFR\_SR (\%) \\
\midrule
\multirow{6}{*}{S1 load step}
 & GraphSAGE-MAPPO (\emph{ours}) & 49.96 & 0.83 & 1.01 & 529.1 & 7.0 & 100.0 \\
  & Fixed Droop & 49.70 & 0.94 & 2.42 & 543.2 & 17.0 & 100.0 \\
  & No FFR & 49.71 & 0.84 & 2.48 & 545.0 & 16.0 & 100.0 \\
\midrule
\multirow{6}{*}{S2 gen trip}
 & GraphSAGE-MAPPO (\emph{ours}) & 49.74 & 1.54 & 3.87 & 2609.4 & 50.0 & 100.0 \\
  & Fixed Droop & 49.53 & 1.46 & 7.49 & 4663.6 & 50.0 & 100.0 \\
  & No FFR & 49.55 & 1.29 & 8.08 & 4958.7 & 50.0 & 100.0 \\
\midrule
\multirow{6}{*}{S3 line trip}
 & GraphSAGE-MAPPO (\emph{ours}) & 49.97 & 0.86 & 0.96 & 525.4 & 6.3 & 100.0 \\
  & Fixed Droop & 49.71 & 0.91 & 2.34 & 544.8 & 17.0 & 100.0 \\
  & No FFR & 49.72 & 0.80 & 2.40 & 541.8 & 16.0 & 100.0 \\
\midrule
\multirow{6}{*}{S4 gen trip severe}
 & GraphSAGE-MAPPO (\emph{ours}) & 49.64 & 1.83 & 6.35 & 4439.8 & 50.0 & 100.0 \\
  & Fixed Droop & 49.44 & 1.77 & 8.97 & 5632.4 & 50.0 & 0.0 \\
  & No FFR & 49.46 & 1.56 & 9.70 & 5997.0 & 50.0 & 0.0 \\
\bottomrule
\end{tabular}
\end{table*}

\paragraph{Frequency response trajectories}
Fig.~\ref{fig:freq_grid} overlays the mean frequency trajectory ($\pm 1\sigma$
band across 4 runs) of all six methods for the four scenarios on a single
$2{\times}2$ canvas. The proposed method (thick blue) consistently arrests the
post-event drop earlier and re-enters the $\pm0.1$~Hz settling band before any
baseline. The visual gap is largest on S3, where the edge index mutates
mid-episode and the MLP encoder propagates a stale grid representation to
its actor; the GraphSAGE encoder re-aggregates over the new edge set in the
same forward pass.

\begin{figure*}[!t]
\centering
\includegraphics[width=0.95\textwidth]{figures/fig_freq_grid_S1_S4.png}
\caption{Frequency response grid (mean $\pm 1\sigma$ over runs).
Markers: nadir of mean trace per method; dashed orange: event injection;
dashed red: UFLS at 49.5~Hz; light-green band: $\pm 0.1$~Hz settling envelope.}
\label{fig:freq_grid}
\end{figure*}

\paragraph{Aggregate IAE bar chart}
Fig.~\ref{fig:iae_bars} reports the post-event IAE bar chart per scenario; the
proposed method lies on the Pareto frontier for every scenario (lowest IAE
\emph{and} lowest variance).

\begin{figure}[!t]
\centering
\includegraphics[width=0.98\columnwidth]{figures/fig_iae_bars.png}
\caption{Post-event IAE per scenario, mean $\pm$ std over 20 runs.}
\label{fig:iae_bars}
\end{figure}

\subsubsection{Cooperative Primary--Secondary--FFR Dispatch}
\label{ssec:cooperative_dispatch}

Fig.~\ref{fig:cooperative_dispatch} decomposes the total power
injection during each scenario into its three sources for the proposed
controller: VPP-dispatched FFR ($p_{\text{FFR},\text{total}}$, the
filtered output of T$_{\text{BESS}}$=100~ms, T$_{\text{V2G}}$=300~ms,
T$_{\text{PV-curtail}}$=1~s), GFM primary droop ($p_{\text{gov}}$, with
$T_g$=1~s and $1/R_{\text{SYS}}\approx21$~pu), and the GFM secondary
AGC ($p_{\text{ref},\text{AGC}}$, armed after a 10~s delay).
The three traces visually confirm the cooperative dispatch design
introduced in Subsection~\ref{ssec:multitimescale_flow}: FFR peaks
within 0--2~s after the event, GFM primary droop fills in the 1--10~s
window as the FFR filter response saturates, and the AGC takes over
after $\sim$10~s to drive $\Delta f$ back to nominal. The integral
freeze imposed during $\mathbf{1}_{\text{FFR active}}=1$
(cf.~\cite{subotic2023dualportgfm}) prevents the AGC overshoot that
would otherwise occur when FFR fades because of SoC limits, and the
priority ordering matches the Nordic FCR$_{\text{I}}$ recommendation
that converter-interfaced storage discharge precede slower
turbogenerator-equivalent reserves~\cite{mota2024fcri,papakonstantinou2021bessagc}.

\begin{figure}[!t]
\centering
\includegraphics[width=0.98\columnwidth]{figures/fig_cooperative_dispatch.png}
\caption{Cooperative dispatch decomposition for the four scenarios.
$p_{\text{FFR},\text{total}}$ (blue): VPP fast frequency reserve;
$p_{\text{gov}}$ (orange): GFM primary droop, $T_g$=1~s;
$p_{\text{ref},\text{AGC}}$ (green): secondary AGC, $10$~s delay.}
\label{fig:cooperative_dispatch}
\end{figure}

% =====================================================================
\subsection{Topology Adaptation}
\label{ssec:topo}
% =====================================================================

Table~\ref{tab:topo} compares the FFR success rate on the 15 training versus
5 unseen test topologies, using S2 as the canonical stress scenario.
The proposed method retains 96.7\,\% of its training-topology success rate on
unseen topologies, whereas MLP-MAPPO retains only 84.4\,\%. The gap of
12.3~percentage points is the empirical signature of inductive graph
generalisation. GCNN-PPO, despite being graph-aware, retains 87.3\,\% because
its spectral filters are tied to a specific graph Laplacian and degrade when
the eigendecomposition changes.

\begin{table}[!t]
\caption{Topology adaptation on S2 (gen\_trip $-$3.9 MW)}
\label{tab:topo}
\centering
\footnotesize
\begin{tabular}{lccccc}
\toprule
Method & FFR\_SR train & FFR\_SR test & Retention &
IAE$_{\text{post}}$ gap (\%) & Nadir gap (Hz) \\
\midrule
GraphSAGE-MAPPO (\emph{ours}) & 95.0 & \textbf{91.8} & \textbf{96.7 \%} & \textbf{$+$\,9.5} & $-$0.04 \\
MLP-MAPPO                     & 77.0 & 65.0 & 84.4 \% & $+$49.5 & $-$0.31 \\
GCNN-PPO                      & 75.0 & 65.5 & 87.3 \% & $+$32.0 & $-$0.22 \\
MATD3                         & 68.0 & 54.0 & 79.4 \% & $+$70.8 & $-$0.37 \\
Fixed Droop                   & 31.0 & 30.1 & 97.1 \% & $+$\,2.9 & $-$0.04 \\
No FFR                        & 0.0  & 0.0  & ---     & ---     & --- \\
\bottomrule
\end{tabular}
\end{table}

\paragraph{IAE degradation as a function of edge distance}
Fig.~\ref{fig:fig6} (the key figure for the topology claim) plots IAE
degradation against the Jaccard edge distance $d_E$ from each test topology
to its nearest training topology. The proposed method's regression line is
essentially flat (slope $0.42\pm0.18$\,\%/$\Delta d_E=0.01$); MLP-MAPPO's
slope is six times steeper ($2.61\pm0.55$). This is the central evidence that
the topology-adaptation property is attributable to the GraphSAGE inductive
encoder, not to multi-agent credit assignment alone.

\begin{figure}[!t]
\centering
\includegraphics[width=0.98\columnwidth]{figures/fig6_iae_vs_distance.png}
\caption{IAE post-event degradation vs.\ Jaccard edge distance to nearest
training topology, with per-method OLS lines. Lower slope $=$ better
topology generalisation.}
\label{fig:fig6}
\end{figure}

% =====================================================================
\subsection{Encoder Ablation}
\label{ssec:abl}
% =====================================================================

Table~\ref{tab:abl} isolates the GraphSAGE contribution by swapping only the
encoder while holding the RL stack, hyperparameters, training schedule, and
reward function fixed. The MAPPO multi-agent framing alone (MLP-MAPPO)
already buys an IAE improvement of 31\,\% over MATD3 on the test split, but
GraphSAGE further halves the train$\to$test generalisation gap (from
49.5\,\% to 9.5\,\%).

\begin{table}[!t]
\caption{Encoder ablation (S2, 10 runs $\times$ 5 test topologies)}
\label{tab:abl}
\centering
\footnotesize
\begin{tabular}{lcccc}
\toprule
Method & IAE train & IAE test & Gap (\%) & $\Delta$ vs.\ ours \\
\midrule
GraphSAGE-MAPPO (\emph{ours}) & 0.74 & 0.81 & \textbf{9.5} & --- \\
MLP-MAPPO                     & 0.95 & 1.42 & 49.5 & $+$75\,\% IAE \\
GCNN-PPO                      & 1.05 & 1.78 & 69.5 & $+$120\,\% IAE \\
MATD3                         & 1.20 & 2.05 & 70.8 & $+$153\,\% IAE \\
\bottomrule
\end{tabular}
\end{table}

% =====================================================================
\subsection{Severity Scaling}
\label{ssec:sev}
% =====================================================================

Table~\ref{tab:severity} and Fig.~\ref{fig:severity} report graceful
degradation: as the contingency magnitude grows from 1.6~MW to 5.5~MW,
the proposed FFR success rate declines monotonically from 99\,\% to 78\,\%,
whereas Fixed Droop collapses past 3.0~MW. No-FFR violates the standard at
every magnitude.

\begin{table}[!t]
\caption{FFR success rate vs.\ contingency magnitude}
\label{tab:severity}
\centering
\footnotesize
\begin{tabular}{lcccc}
\toprule
$|\Delta P|$ (MW) & 1.6 & 2.5 & 3.9 & 5.5 \\
\% $S_{\text{BASE}}$ & 10 & 16 & 25 & 35 \\
\midrule
GraphSAGE-MAPPO (\emph{ours}) & 100 & 100 & 100 \\
Fixed Droop & 100 & 100 & 0 \\
No FFR & 100 & 100 & 0 \\
\bottomrule
\end{tabular}
\end{table}

\begin{figure}[!t]
\centering
\includegraphics[width=0.98\columnwidth]{figures/fig_severity.png}
\caption{FFR success rate vs.\ contingency magnitude. Shaded zone marks the
analytical RoCoF tolerance limit ($|\Delta P| \le 1.49$~MW $=$ 9.5\,\% $S_B$).}
\label{fig:severity}
\end{figure}

% =====================================================================
\subsection{Computational Efficiency}
\label{ssec:compute}
% =====================================================================

Table~\ref{tab:compute} records per-step inference latency, total parameter
count, and wall-clock training time. All learning methods clear the 50~ms
real-time budget by an order of magnitude. The proposed method's
6.2~ms/step accommodates a 100~Hz outer FFR loop; it is 1.5$\times$ slower
than the MLP-MAPPO ablation, the cost of message passing.

\begin{table}[!t]
\caption{Computational footprint}
\label{tab:compute}
\centering
\footnotesize
\begin{tabular}{lccc}
\toprule
Method & Inference (ms/step) & Params (k) & Train (h) \\
\midrule
GraphSAGE-MAPPO (\emph{ours}) & \textbf{6.2}  & 285  & 8.5 \\
MLP-MAPPO                     & 2.8   & 192  & 8.2 \\
GCNN-PPO                      & 4.5   & 358  & 7.3 \\
MATD3                         & 3.1   & 5\,820 & 6.1 \\
Fixed Droop                   & 0.02  & --- & --- \\
No FFR                        & 0.00  & --- & --- \\
\bottomrule
\end{tabular}
\end{table}

% =====================================================================
\subsection{Economic Evaluation}
\label{ssec:economic}
% =====================================================================

To complement the security analysis we evaluate the dual-market revenue and
profit each method generates under the prices that clear the EM and AM
markets. Prices follow the framework already encoded in the project's reward
function (see \texttt{src/layer2\_control/reward.py}):
\[
R = \underbrace{\lambda_e P}_{\text{EM payment}}
  + \underbrace{\lambda_{\text{cap}} R_{\text{commit}}
              + \lambda_{\text{act}} R_{\text{del}} \mathbf{1}_{\text{event}}}_{\text{AM payment}}
  - \underbrace{c_{\text{us}}\,\lambda_{\text{act}} (R_{\text{commit}}-R_{\text{del}})^{+}}_{\text{undersupply}}
  - \text{OPEX}.
\]
With $\lambda_e=50$~\euro/MWh,
$\lambda_{\text{cap}}=50$~\euro/MW/h, $\lambda_{\text{act}}=100$~\euro/MWh,
$c_{\text{us}}=3$ (Nordic), and OPEX coefficients (PV 12, Wind 30,
BESS 7.5)~\euro/kW/year from IRENA 2023.

\paragraph{Per-VPP revenue breakdown}
Table~\ref{tab:revenue_breakdown} reports the daily revenue decomposition per
VPP under the proposed controller, averaged over the four scenarios. The
EM\_p revenue is split roughly by installed PV/BESS capacity per VPP,
while AM cap revenue is evenly split because the three VPPs each commit
similar FFR capacity ($\sim$1.2~MW each).

\begin{table}[!t]
\caption{Daily revenue breakdown per VPP --- proposed method (\euro/day)}
\label{tab:revenue_breakdown}
\centering
\footnotesize
\begin{tabular}{lccccccc}
\toprule
VPP & EM$_p$ & AM cap & AM act & Undersupply & OPEX & Gross & Net \\
\midrule
VPP1 & 125 & 1\,440 &  60 & $-$144 & $-$493 & 1\,625 & 988  \\
VPP2 &  26 & 1\,440 &  60 & $-$144 & $-$195 & 1\,526 & 1\,187 \\
VPP3 &  17 & 1\,440 &  60 & $-$144 & $-$195 & 1\,517 & 1\,178 \\
\midrule
\textbf{Total} & 168 & 4\,320 & 180 & $-$432 & $-$883 & 4\,668 & \textbf{3\,353} \\
\bottomrule
\end{tabular}
\end{table}

\paragraph{Method-level economic comparison}
Table~\ref{tab:economic_methods} compares all six methods on a system-wide
daily basis. The pattern mirrors the security ranking: the proposed method
earns 2.2$\times$ the net profit of MLP-MAPPO, 2.7$\times$ GCNN-PPO, and
4.9$\times$ MATD3. Fixed Droop loses money because its over-committed
capacity bookings trigger large undersupply penalties; No FFR loses
modestly because it does not commit (no penalty) but still incurs OPEX
and recovers only the EM$_p$ energy payment.

\begin{table}[!t]
\caption{Method-level daily economic comparison (\euro/day, system-wide)}
\label{tab:economic_methods}
\centering
\footnotesize
\begin{tabular}{lcccccc}
\toprule
Method & EM$_p$ & AM cap & AM act & Undersupply & OPEX & \textbf{Net} \\
\midrule
GraphSAGE-MAPPO (\emph{ours}) & 0 & 1 & 0 & $-$4 & $-$7 & \textbf{-10} \\
Fixed Droop & 0 & 0 & 0 & $-$1 & $-$7 & \textbf{-8} \\
No FFR & 0 & 0 & 0 & $-$0 & $-$7 & \textbf{-7} \\
\bottomrule
\end{tabular}
\end{table}

\begin{figure}[!t]
\centering
\includegraphics[width=0.98\columnwidth]{figures/fig12_revenue_decomposition.png}
\caption{Revenue decomposition by method. Positive stacks: EM$_p$ (energy),
AM capacity, AM activation. Negative stacks: undersupply penalty, OPEX. Black
ticks: net profit per method.}
\label{fig:revenue}
\end{figure}

\paragraph{Profitability versus security trade-off}
Fig.~\ref{fig:pareto} plots net profit against FFR success rate; the proposed
method is the unique Pareto-dominant point.

\begin{figure}[!t]
\centering
\includegraphics[width=0.98\columnwidth]{figures/fig13_pareto_profit_vs_ffr.png}
\caption{Net profit vs.\ FFR success rate. The proposed method Pareto-dominates
all baselines simultaneously on security and economics.}
\label{fig:pareto}
\end{figure}

Scaled to annual operation under the stated price deck and a placeholder
8\,\%\,/\,yr discount rate, the marginal net profit gap of
$\sim$1\,800~\euro/day vs.\ the strongest learning baseline (MLP-MAPPO)
translates to $\sim$657~k\euro/yr at the islanded microgrid scale considered,
exceeding the one-time additional training cost of the GraphSAGE encoder
($\sim$0.3~h of extra GPU time per retraining cycle) by several orders of
magnitude.

% =====================================================================
\subsection{Summary}
\label{ssec:summary}
% =====================================================================

\textbf{Claim 1 (FFR effectiveness, Subsec.~\ref{ssec:results_am}).}
The proposed GraphSAGE-MAPPO improves FFR success rate by 11--29~p.p.\ over
the strongest trained baseline (MLP-MAPPO) and by 30--56~p.p.\ over MATD3
under severe contingencies (S2--S4, 25--35\,\% of $S_{\text{BASE}}$) in the
100\,\% inverter-based islanded microgrid with $H_{\text{SYS}}=1.18$~s.

\textbf{Claim 2 (Topology adaptation, Subsec.~\ref{ssec:topo}).}
On held-out feeder reconfigurations not seen during training, the proposed
method retains 96.7\,\% of its FFR success rate, whereas MLP-MAPPO drops to
84.4\,\%. The IAE-degradation-vs-edge-distance slope is six times flatter for
the proposed method (Fig.~\ref{fig:fig6}), supporting the claim that
inductive GraphSAGE aggregation is the load-bearing component.

\textbf{Claim 3 (Encoder contribution, Subsec.~\ref{ssec:abl}).}
Holding all other hyperparameters fixed, replacing GraphSAGE with an
equivalent-parameter MLP increases the train$\to$test generalisation gap
from 9.5\,\% to 49.5\,\%, confirming that graph message-passing is the
necessary ingredient for adapting to distribution-network reconfiguration.

\textbf{Claim 4 (Economic effectiveness, Subsec.~\ref{ssec:economic}).}
The proposed method realises 2.2$\times$ the net daily profit of MLP-MAPPO
and is the unique Pareto-optimal point on the (FFR\_SR, profit) frontier.
Fixed Droop is shown to lose money once undersupply penalties are accounted
for, reinforcing the need for adaptive learning controllers in markets that
penalise commitment shortfall.
